\documentclass[12pt,letterpaper]{article}
\usepackage[utf8]{inputenc}
\usepackage[T1]{fontenc}
\usepackage{times}
\usepackage[margin=0.5in,top=0.4in,bottom=0.4in]{geometry}
\usepackage{amsmath,amssymb,mathrsfs,braket}
\DeclareMathOperator{\tr}{tr}
\DeclareMathOperator{\rank}{rank}
\DeclareMathOperator{\diag}{diag}
\usepackage{hyperref}
\usepackage{xcolor}
\usepackage{titlesec}
\hypersetup{colorlinks=true,urlcolor=blue!60!black}
\setlength{\parindent}{0pt}
\setlength{\parskip}{3pt}
\titleformat{name=\section,numberless}{\large\bfseries\color{blue!60!black}}{}{0pt}{}[\vspace{-2pt}\rule{\linewidth}{0.6pt}]
\titlespacing{\section}{0pt}{8pt}{4pt}
\pagestyle{plain}
\begin{document}
{\footnotesize\textbf{Quantum Computing Digest --- 2026-11-05} \hfill 7 papers}

{\scriptsize Full author abstracts}

\section*{Machine Learning for Quantum Control and Pulse Optimization}

{\small\textbf{Efficient training of a 128-modes 3D continuously coupled photonic chip via limited quantum resources}} {\scriptsize[\href{https://arxiv.org/abs/2611.05681}{arxiv}]}\nopagebreak\\
{\scriptsize Denis Stanev, Simone Di Micco, Francesco Hoch, Taira Giordani, Nicol\`o Spagnolo, Marco Gardina, Andrea Crespi, Francesco Ceccarelli, Roberto Osellame, Fabio Sciarrino}\nopagebreak\\[2pt]
{\footnotesize Integrated photonics provides a promising route towards the development of large-scale, stable and low-loss quantum hardware. Existing implementations of universal quantum photonic chips generally rely on architectures consisting of localised interactions between the modes that can be programmed with known algorithms to implement specific operations. Recent advances have shown that interferometers based on 3D continuously coupled waveguide arrays may allow to significantly scale up the number of modes with reduced overhead. However, for this class of devices, optimal strategies to set the controllable parameters and obtain the desired transformation have not been demonstrated generally. In this work, we report the use of a machine learning algorithm driven by quantum data to program the operation of a reconfigurable three-dimensional continuously coupled photonic device comprising 128 modes. We show experimentally the capability of reaching a target unitary matrix, among those achievable by the device, through a limited number of one- and two-photon measurements. Our results demonstrate that this approach results to be effective in the presence of noise and with a limited amount of resources, also allowing to partially compensate imperfections in the experimental setup. This method can thus represent a viable route towards programming large scale integrated photonic devices, unlocking the possibility of controlling photonic systems for which an analytical decomposition is not trivial.\par}
\medskip


\section*{Machine Learning for Error Correction, Mitigation and Fault Tolerance}

{\small\textbf{Discovering correlated error mechanisms: data-driven augmentation of detector error models improving decoding}} {\scriptsize[\href{https://arxiv.org/abs/2611.06605}{arxiv}]}\nopagebreak\\
{\scriptsize Zohar Schwartzman-Nowik, Matthew P. Stafford, Boris M. Varbanov, Joan Camps}\nopagebreak\\[2pt]
{\footnotesize Accurate decoder priors are essential for high-performance quantum error correction. Theoretical noise models inevitably omit physically relevant error mechanisms, most notably correlated errors, which appear consistently in experimental error budgets yet are absent from standard models. We introduce a data-driven method for augmenting detector error models (DEMs) by discovering entirely new error mechanisms directly from experimental detector-firing statistics, going beyond most existing approaches that merely refine the probabilities of already-known mechanisms. Our method generates candidates by combining pairs of existing mechanisms as building blocks, then filters and validates them using pairwise and higher-order detector correlations. Crucially, this construction provides each newly discovered mechanism with a well-defined logical effect derived from the constituent error mechanisms, enabling its direct use in decoding. Furthermore, each mechanism is accompanied by a natural decomposition into lower-order components, preserving compatibility with the original decoder. This is particularly valuable in settings that rely on correlated graph-based decoders such as correlated matching. We evaluate the method on surface-code simulations with correlated noise and correlated-matching decoding, demonstrating that, in a realistic near-term noise scenario, DEM augmentation increases the Lambda parameter by 20\% across code distances 7-15 and reduces the logical error rate by 75\% at code distance 15 relative to the unaugmented DEM. On experimental data from Google's superconducting device, the augmented DEM increases the Lambda parameter by 9\% across code distances 5-7 and yields a reduction of 14\% in logical error rate over the baseline for distance 7. To our knowledge, this is the first data-driven DEM augmentation method that produces decomposed hyperedges with assigned logical effects.\par}
\medskip

{\small\textbf{Trinetra.Q: Statistical and Behavioral Detection of Attacks on the Quantum-Error-Correction Layer and Classical Control Stack of Quantum Computing Systems}} {\scriptsize[\href{https://arxiv.org/abs/2611.05722}{arxiv}]}\nopagebreak\\
{\scriptsize Garrv Sipani}\nopagebreak\\[2pt]
{\footnotesize Most quantum-computer security work targets physical-layer attacks on qubits, while two surfaces remain inadequately defended: the quantum-error-correction (QEC) syndrome-and-decoder pipeline, and the classical control and software stack that runs it. Tampering on these surfaces can corrupt a computation while the physical error rate appears normal, and can stay within a single-circuit integrity tolerance. We present Trinetra.Q, a detection layer of cooperating detectors, each monitoring a statistical or behavioral property of quantum-domain data against an envelope calibrated from honest operation, and raising an indication only on a statistically significant deviation. The detectors cover (A) QEC integrity over per-stabilizer firing rates, syndrome defect-count, and repeated-syndrome rate; (B) decoder confidence over the margin between minimum-weight error explanations in distinct logical classes; (C) behavioral output divergence, flagging crosstalk-class and novel attacks by their effect on a victim's output relative to a matched control; (D) control-stack tampering via known-answer canary circuits, built to be indistinguishable to a resident implant, with a cross-job consistency statistic; and (E) compilation/execution integrity. Each gates on significance, so honest operation raises no alarm. We evaluate in faithful simulation and on three production IBM Heron processors, and bound the physical attack surface with a systematic negative map: most physical-layer attacks return statistically null, the exception an implemented, measured within-job feed-forward covert-channel primitive. Every detector is also evaluated against an adaptive attacker with full knowledge of it. To our knowledge this is the first unified statistical-and-behavioral detection layer for the QEC layer and classical control stack of a quantum computing system, demonstrated on production hardware.\par}
\medskip


\section*{Notable Quantum Computing Research}

{\small\textbf{Unlocking a fast adiabatic CZ gate and exact residual $ZZ$ cancellation between fixed-frequency transmons using a floating tunable coupler}} {\scriptsize[\href{https://arxiv.org/abs/2604.05048}{arxiv}]}\nopagebreak\\
{\scriptsize Angela Q. Chen, Sarah Strong, Xian Wu, Stefano Poletto}\nopagebreak\\[2pt]
{\footnotesize Tunable couplers in superconducting qubit architectures enable strong qubit-qubit interactions for two-qubit gates while suppressing unwanted coupling during single-qubit operations. However, achieving low error rates for fast two-qubit gates remains challenging, as suppressing leakage and non-adiabatic errors typically requires specialized qubit, coupler, or pulse designs, often at the expense of an idling $ZZ=0$ condition. In this work, we demonstrate that a symmetric floating tunable coupler provides a natural platform for fast, high-fidelity adiabatic controlled-Z (CZ) gates. Its favorable energy-level structure eliminates the conventional trade-off between rapid conditional-phase accumulation and adiabatic evolution while preserving exact cancellation of residual $ZZ$ interaction at idling. This architecture exhibits intrinsic robustness to non-adiabatic transitions, even under simple flux modulation waveforms. To push performance at short gate durations, where maintaining adiabaticity becomes more challenging despite the favorable level structure, we introduce pulse-shaping techniques based on the instantaneous adiabatic factor that further suppress non-adiabatic errors. We experimentally realize a 24 ns adiabatic CZ gate with fidelity exceeding 99.9\% and stable operation over several hours with a time-averaged fidelity of 99.878\%\par}
\medskip

{\small\textbf{Flux Crosstalk Calibration from Zero Knowledge to High Precision in Fully Tunable Superconducting Circuits}} {\scriptsize[\href{https://arxiv.org/abs/2611.06656}{arxiv}]}\nopagebreak\\
{\scriptsize Y. Gao, A. Galicia, J. R. Guimar\~aes, Y. Haddad, Y. Liu, J. Cereijo, D. A. Volkov, H. Bhardwaj, M. Neis, M. Jerger, P. A. Bushev, R. Barends}\nopagebreak\\[2pt]
{\footnotesize Precise and independent flux control is essential for scalable superconducting processors built from frequency-tunable qubits and couplers, but flux crosstalk undermines it. Conventional calibration methods rely either on dedicated readout channels, which tunable couplers typically lack, or on readout schemes requiring extensive prior spectroscopy of the qubit and coupler. We demonstrate two independent methods that characterize coupler flux crosstalk without directly driving or measuring the couplers. The first method is zero-knowledge characterization, which exploits intrinsic coupler-resonator coupling to yield rapid and accurate estimates with no prior qubit or coupler spectroscopy, accelerating large scale bring-up. The second method is symmetric branch cancellation, which isolates crosstalk from other dispersive effects allowing for quantifying with high precision. Combined with an improved iterative procedure, these methods suppress residual coupler crosstalk below $5\times10^{-5}$, providing a complete route from zero-knowledge initialization to high-precision crosstalk characterization in fully tunable architectures.\par}
\medskip

{\small\textbf{Resetting a Superconducting Qubit with a Compact Colder Resistive Dissipator below the Qubit Frequency}} {\scriptsize[\href{https://arxiv.org/abs/2611.06216}{arxiv}]}\nopagebreak\\
{\scriptsize Jiajia Tian, Meiling Li, Weijie Sun, Yanjun Ma, Zhenyu Mi, Fei Yan, Huikai Xu, Pei Liu, Guangming Xue, Haifeng Yu}\nopagebreak\\[2pt]
{\footnotesize Efficient qubit reset has become a key requirement for large-scale superconducting quantum processors, as initialization fidelity increasingly lags behind gate performance and quantum error correction demands rapid, high-fidelity state re-initialization. This work proposes a compact on-chip dissipator consisting of two transmon dissipation modes and a thin-film resistor, and experimentally demonstrates an unconditional reset protocol. With no additional control lines and negligible impact on the qubit $T_1$, the effective temperature of the transmon qubit is reduced from $49$\textasciitilde{}mK to $35$\textasciitilde{}mK. Quantum process tomography is employed to benchmark the reset process, achieving a process error of $0.0028$. This reset protocol is robust to pulse parameter deviations, thereby relaxing calibration requirements. The dissipator architecture provides a practical, calibration-friendly reset solution for large-scale superconducting quantum processors.\par}
\medskip

{\small\textbf{Solving a non-abelian hidden subgroup problem on a trapped-ion quantum computer: the power of the measurement basis}} {\scriptsize[\href{https://arxiv.org/abs/2611.06716}{arxiv}]}\nopagebreak\\
{\scriptsize Oliver O'Brien}\nopagebreak\\[2pt]
{\footnotesize The hidden subgroup problem (HSP) unifies many of the best-known quantum speed-ups. For abelian groups it is solved efficiently by weak Fourier sampling, but for non-abelian groups, which include the cases relevant to graph isomorphism and lattice-based cryptography, the right measurement basis is crucial and often unknown. We experimentally study the HSP over the $q$-hedral groups $\mathbb{Z}_q \ltimes \mathbb{Z}_p$, which interpolate between the dihedral and affine groups and mark a boundary between instances that are efficiently solvable and those that are only information-theoretically solvable. On the IonQ Forte trapped-ion quantum computer, with $p \le 7$, we implement weak Fourier sampling, strong Fourier sampling in an adapted basis, strong Fourier sampling of adjacent pairs, and forgetful abelian sampling. We demonstrate that the theoretical predictions of Moore, Rockmore, Russell and Schulman survive on a noisy device: weak sampling reveals nothing about the hidden subgroup, the adapted basis succeeds for large $q$, adjacent pairs succeed for small $q$, and forgetful sampling saturates below unit success. For $p=7$, $q=6$, the hardware exceeds the success probability of an exhaustively optimised classical algorithm for up to three oracle queries.\par}
\medskip


\end{document}